\chapter{Introduzione}
\label{cha:Introduzione}

La sincronizzazione temporale è fondamentale nei sistemi distribuiti e nelle
reti di calcolatori, consetne a nodi differenti di avere un sistema di riferimento
temporale comune sufficientemente coerente. In assenza di un riferimento
temporale condiviso, diventa più complesso correlare eventi prodotti da macchine diverse,
ordinare correttamente operazioni distribuite, analizzare log, coordinare
servizi e valutare il comportamento complessivo di un'infrastruttura di rete.

Nei sistemi distribuiti e nelle reti di calcolatori, rispetto ad ognuno di essi vi
è una memoria del proprio clock locale, producendo infomazioni temporali
associate ad ogni file, log, transazione, messaggi e processi automatici. I
timestamp, rappresentativi appunto dell'istante temporale in cui è avventuo un qualche
cambiamento rispetto all'informazione, fungono da riferimento tra macchine
differenti, venendo trattati come se fossero direttamente confrontabili tra macchine
differenti, tale assunzione però può essere valida solamente se i clock locali
vengono tenuti aggiornati rispetto ad un riferimetno comune.

La difficolta alla base è più complessa di quanto può sembrare. I clock locali non
sono strumenti perfetti, e in quanto tale la loro stabilità e correttezza nel
tener conto della miusra del tempo può essere incerta. Gli oscillatori hardware possono
accumulare deriva nel tempo, generando scostametni progressivi rispetto ad un
tempo di riferimento. La sincronizzazione, quindi, non è solamento uno strumento di
inizializzazione per l'orologio di una macchina, ma è uno strumento necessario
per compensare in modo continuo la deriva tra orologia diversi.

L'importanza della sincronizzazione temporale è centrale in una molitudine di
ambiti applicativi. Nei sistemi di logging, auditing e monitoraggio, timestamp coerenti
consentono di ordinare eventi prodotti da host diversi e di ricostruire correttamente
il comportamento della rete. Nella diagnosi dei guasti, la possibilità di correlare
temporalmente eventi registrati da server, router e altri dispositivi riducendo
drasticamente l'ambiguità nell'identificazione della causa del problema. Servizi come
quelli di directory, autenticazione, operazioni schedulate, ambienti
virtualizzati, posta elettronica e sistemi transazionali dipendono, in
condizioni differenti, da una base di riferimento comune.

Una sbagliata sincornizzazione può essere causa di problematicità, la quale può
avere dimensioni differenti rispetto all'ambito di riferimento in cui si opera. Vi
sono sistemi che sono dotati di tolleranze agli errori, arrivando a
discostamenti fino all'ordine di millisecondi, tuttavia trascuravili per attività
come quelle di logging, monitoraggio o correlazione approssiamta degli eventi. Vi
sono poi scenari, come sistemi industriali, telecomunicazini, applicazioni real-time
o infrastrutture finanziarie che richiedono alta precisione e una maggiore
stabilità, andando a ricercare i sistemi più sofisticati che possano mitigare a tali
problematicità. In tali casi, ritardi variabili, jitter, perdita o riordinametno dei
pacchetti possono compromettere la qualità della sincronizzazione e rendere necessario
valutare attentamente il comportamento dei diversi protocolli in condizioni di rete
non ideali.